\chapter{The Reproduction Harness}
\label{app:harness}

\keybox{what this appendix is for}{%
Nothing here is a result. This is the machinery that produced the results: what
runs where, what is checked before anything is measured, how a measurement gets
into the record and how it gets out again, and what the compute allowance
allowed. A reader who takes Chapter~\ref{ch:results} at face value can skip it.
A reader who wants to know whether the numbers can be trusted --- or who wants
to re-run them --- will find the answer here.}

\section{Where The Work Runs}

There is no local GPU. Every training run in this thesis was executed on Kaggle's
free tier: two NVIDIA T4s, $15\,360$\,MiB each, in sessions capped at twelve
hours, against a weekly allowance of $30$ GPU-hours per account. That shaped the
design more than any other single fact, and it is stated first because several
decisions later in this appendix look arbitrary until you know it.

A \emph{rung} is one such session. It clones the repository at a named commit,
builds the CUDA rasteriser from source, converts the dataset, trains, renders,
evaluates, writes one CSV of results and one JSON summary, and prunes everything
else so that the session's output fits in what the platform will hand back. The
scripts are \texttt{kaggle/*.py}; the launcher that pushes them and collects
their output is \texttt{tools/kaggle\_push.py}.

\section{Pre-Flight Assertions}
\label{app:preflight}

An instrument has to be checked before it is used. Twelve assertions run before
the first training step and abort the session if any fails. Each one exists
because the thing it checks can fail silently, and several of them have.

\input{generated/table-preflight}

Three of them:

\begin{itemize}[leftmargin=1.4em,itemsep=3pt]
  \item \textbf{The accelerator check.} Kaggle allocates whichever GPU is free.
        3DGS needs compute capability $7.0$ and a P100 is $6.0$, so a session
        landing on one would fail forty minutes in with a message that looks like
        a code fault. It now aborts in thirty seconds with a marker the launcher
        recognises and relaunches on.
  \item \textbf{The arm-separation diff check.} The 3DGS arm must differ from the
  Mip-Splatting arm by exactly the intended change and nothing else. The check enumerates the files that
        change and refuses anything outside the list. It has fired twice in
        anger --- once on a runtime state file that a careless \texttt{git add}
        had committed, and once when the switch of Appendix~\ref{app:impl}
        reached further than its author believed.
  \item \textbf{The LPIPS backbone.} LPIPS defaults to AlexNet; both papers
        report VGG. The two differ by more than the effects this thesis
        measures. The backbone is pinned, asserted, and written into every row
        of the record as a column, so a row cannot be compared against a
        published figure without the comparison being checkable.
\end{itemize}

\section{The Record}
\label{app:record}

Every measurement lands in one file, \texttt{results/runs.csv}, one row per
(arm, scene, test scale). The file is \textbf{append-only}. A configuration that
is re-measured does not overwrite its predecessor; the predecessor stays and is
marked, in its own \texttt{notes} column, with the name of the stage that
replaced it. Nothing is ever deleted.

That discipline exists because a deleted row is a measurement that cannot be
re-examined, and this project has needed to re-examine several. It also has a
failure mode, which Appendix~\ref{app:defects} describes: a row that is kept but
not marked is a row that is silently \emph{used}. The marking is the only thing
standing between an honest archive and a pooled average, and the marking is now checked --- the generator refuses to emit any table whose
cells draw on more than one build.

\paragraph{From The Record To The Document.} No measured number in this thesis is
typed by hand. \texttt{tools/make\_tex.py} reads \texttt{results/runs.csv} and
writes every table and every inline figure as a LaTeX macro; an unmeasured quantity expands to a visible marker that a reader cannot mistake
for a result. The same loader feeds the figures and the slide deck, so the three cannot
disagree --- and a test asserts that they do not, because for a while they did.

\section{The Apparatus}
\label{app:artefacts}

\begin{table}[htbp]
  \centering
  \caption{What exists, and what each artefact is for.}
  \small
  \begin{tabularx}{\linewidth}{L{0.70}L{1.30}}
    \toprule
    artefact & purpose \\
    \midrule
    \texttt{results/runs.csv} & Append-only. One row per (arm, scene, test-scale, seed), carrying both repository commits, LPIPS backbone, resolution flag, peak VRAM, primitive count, wall-clock and the claimed acceptance level. Every table and figure regenerates from it. \\
    \texttt{tools/make\_tex.py} & Writes every measured table and every inline number in this document. A table with no measured rows is emitted as an explicit placeholder carrying the published targets and nothing in the measured columns; it is never quietly filled from the published column. \\
    \texttt{tools/make\_figures.py} & Generates every figure. In \texttt{--mode measured} it redraws from the CSV with the published curve retained as a faint reference, and refuses to run if the CSV has no usable rows. \\
    \texttt{tools/split\_by\_scale.py} & Recovers the per-scale columns that the shipped evaluation pipeline blends into a single average. Without it there is no multi-scale table. \\
    \texttt{tools/selftest\_pipeline.py} & Runs the whole reporting chain against a fixture with planted, known-in-advance answers, on a machine with no GPU. It is what caught the two silent defects described below. \\
    \texttt{kaggle/*.py} & The rung kernels, each printing every assertion it makes and ending in a single machine-readable summary line. \\
    \bottomrule
  \end{tabularx}
\end{table}

\paragraph{Two Defects The Self-Test Caught.} The
figure generator's ``measured'' mode was drawing the \emph{published} curve under
a measured filename, because \texttt{runs.csv} stores the method as
\texttt{3dgs} while the figures key off \texttt{3DGS}: every lookup missed, the
non-empty guard still passed, and the output was a plausible-looking lie. Second,
the per-scale splitter assumed a JSON nesting that \texttt{metrics.py} does not
write, which would have failed only after both arms had trained. Neither would
have raised an error where it happened. They are recorded here because a
reproducibility claim that has never been tested against a fixture with known
answers is an assertion, not a property.

\section{Cost Per Rung}
\label{app:ladder}

\input{generated/table-budget}

Three entries need a sentence each, because they carry the project's real costs.

\begin{itemize}[leftmargin=1.4em,itemsep=3pt]
  \item \texttt{r3\_real} is the largest single session and did not complete. It
        ran the four indoor Mip-NeRF 360 scenes and was still going when the
        twelve-hour cap arrived. What it produced is kept and reported as a
        four-scene subset, clearly labelled; what it did not produce is not
        estimated.
  \item The 3DGS-arm re-runs after the opacity-compensation fix --- \texttt{c2\_armb} and its continuations
        --- are pure re-measurement cost, incurred because the arm had been built
        wrong. They are the price of the defect in Appendix~\ref{app:defects},
        and they are itemised so that the price is visible.
  \item The stress suite is cheap by comparison. Restricting the training cameras
        does not make training longer, and the sixth protocol --- the one that
        produced the only positive result in this thesis --- trains on three
        cameras and is the fastest run in the table.
\end{itemize}

\paragraph{The Allowance Was Exhausted, Twice.} The first time stalled the
programme for three days and is the reason a second account exists. The binding
constraint on this project has been compute, not code, throughout; the ladder is
ordered so that the cheapest falsifications come first, and the two that could
have ended the project --- does the Nyquist floor bind everywhere, and is the
capture conditioning really anisotropic --- were run before a line of method code
was written.
